\subsection{Replies to comments received from Nakadaira-san - May 17, 2013}

\begin{enumerate}
%........................................................................................................
\vspace{5 mm}
\item
{\em
  I have one comment on the updated part.

  P.49 L852: "The constant $\Delta\chi^2$ method ..." 

  You have described the fitter performance at Appendix H.
  In page 329, the size of the bias for $sin^2(2 \theta_{23})$ is 0.025.
  It is the half of the obtaind 95\% region.
  And the width of the pull of $sin^2(2 \theta_{23})$ is 0.9 and smaller than 1. 

  So, it is difficult to judge the 95\% region obrained from the constant $\Delta\chi^2$ method is reasonable, I think.
  We can not avoid to do F-C method to obtain the reliable contour for $sin^2(2 \theta_{23})$.

  On the other hand, there is no bias for $|\Delta m^2_{23}|$. So, I think that the 1 sigma region for $|\Delta m^2_{23}|$ is reasonable, I think.

  P.49 L852: " .... there is a problem in using the Feldman-Cousins
  method. There is no obvious way of choosing the true
  value(s) of $|\Delta m^2_{23}|$ as inputs to the toy MC
  experiments to find the critical values of $\Delta\chi^2$ for
  the 1D confidence interval in $sin^2(2 \theta_{23})$."

  I think that you can do the 1D F-C method for $sin^2(2 \theta_{23})$ by assuming the best fit value of $|\Delta m^2_{32}|$ for toy MC input.

  Fortunately, the obtained 2D allowed region is not tilted, and it means that there is no correlation between $sin^2(2 \theta_{23})$ and
  $|\Delta m^2_{32}|$ estimation. So, it is expeceted that the results of 1D
  F-C method for $sin^2(2 \theta_{23})$ is not depends on $|\Delta m^2_{32}|$ assumption very much.

  I think that the allowed region the 1D F-C method is more reliable than the region obtained by constant $\Delta\chi^2$ method,
  even 1D F-C method is assuming some $|\Delta m^2_{32}|$ value.
}\\
\\
This reply is also addressed to Anselmo, Kate, Boris and Neil, who requested clarification on the same point.

We agree that Feldman-Cousins is better than the constant-$\Delta \chi^{2}$ method if the best fit is near a physical boundary. 
However there is a problem in the practical implementation of Feldman-Cousins to obtain 1D confidence intervals: there is no 
standard procedure for choosing the true input value(s) of $|\Delta m^{2}_{32}|$ to the toy experiments when obtaining the 
critical values of $\Delta \chi^{2}$ for the interval in $\sin^{2} 2 \theta_{23}$ (and vice versa). Your suggestion of using 
the T2K best-fit values appears at first sight to be the obvious choice, but there are two problems with it. First, a 1D interval 
in $\sin^{2} 2 \theta_{23}$ should be constructed so that the true unknown value is contained in 90\% of all confidence intervals, 
irrespective of the true value of $|\Delta m^{2}_{32}|$. If the value of $|\Delta m^{2}_{32}|$ is fixed, the uncertainty in that 
parameter is not taken into account, and this would give an incorrect confidence interval in $\sin^{2} 2 \theta_{23}$. Second, the 
1D confidence interval in $|\Delta m^{2}_{32}|$ definitely depends on the choice of the true input value of $\sin^{2} 2 \theta_{23}$ 
to the toy experiments. A choice of $\sin^{2} 2 \theta_{23}$ = 1.0 would give considerably greater pile-up at the physical boundary 
when fitting the toy experiments than $\sin^{2} 2 \theta_{23}$ = 0.9, and this would result in a different confidence interval in 
$|\Delta m^{2}_{32}|$.

For these reasons, it is not straightforward to choose the true input values of the ``other'' parameter. Two members of the VaLOR 
group used to work on MINOS, and we know that MINOS came across this same problem. The collaboration spent several weeks trying to 
resolve it but was unable to find a satisfactory solution.
\end{enumerate}
